\section{The mass in the deepest part of the cusp}
\label{dc-section}

We now prove the remaining estimate \eqref{amb-deep-input}.
The argument uses the zero constant term of a discrete-series
automorphic form, a small fundamental domain for one unipotent
radical, and an explicit local Sobolev inequality.  We import
the classical Euclidean transference theorem of Banaszczyk,
the cuspidality theorem for discrete-series matrix coefficients,
and the usual analytic continuation and functional equation
of a finite-order Hecke character.  These imports are stated
precisely where they are used.  The quantitative estimates
and their application to the actual kernel are proved below.

\subsection{Statement and an absolute quotient-volume bound}

Retain the notation of Sections~\ref{amb-section},
\ref{ak-section}, \ref{rv-section}, and \ref{nu-section}.
In particular, \(E/F\) is CM and unramified at every finite
place, \([F:\mathbb Q]=f\), and the Hermitian lattice is
self-dual.  Write
\[
 N=n^2f,\qquad X=nfr_Egb,\qquad \tau_0=X^{-11q}.
\]
Let \(K(x,y)\) be the orthogonal-projection kernel constructed
in Section~\ref{ak-section}, of finite rank \(D\).

\begin{theorem}\label{dc-main}
Suppose \(\log n\ge20\), \(1\le q\), \(q^2\le n\),
\(b\ge2\), \(g\ge1\), and
\(\lvert\lambda_{v,i}\rvert\le b^2\).
Suppose also that the hypotheses of
Section~\ref{ak-section} hold, so its absolutely locally
uniformly convergent projection kernel and finite-rank
identity are available.  Then
\begin{equation}\label{dc-main-bound}
        \int_{\lambda_x<\tau_0}K(x,x)\,dx
                                    \le e^{-N}D.
\end{equation}
No upper bound on \(f,r_E,g\), or \(\log b\) is needed
in the estimates of this section.
\end{theorem}

We may normalize finite Haar measure by
\(\operatorname{vol}(K_f)=1\), and archimedean Haar measure
by the Hilbert--Schmidt density on each real Lie algebra.
Changing Haar measure multiplies an integral kernel and
its diagonal by the reciprocal factor and leaves its
operator, rank, and integrated diagonal unchanged.
Thus this normalization does not change
\eqref{dc-main-bound}.

Let \(\chi\) be the quadratic Hecke character of \(E/F\),
and put \(d_F=|D_F|=r_E^f\).  It has finite conductor one
and is odd at every real place.

\begin{lemma}\label{dc-volume}
In the preceding normalization,
\begin{equation}\label{dc-volume-bound}
 L_F(1,\chi)\le3^f d_F^{1/4},\qquad
 \operatorname{vol}(\mathcal X)\le X^N,\qquad
 \operatorname{vol}(K_f(g))^{-1}\le g^{2N}.
\end{equation}
\end{lemma}
\begin{proof}
We use the classical Hecke theorem in the following form:
for a nontrivial primitive finite-order character, the
completed \(L\)-function is entire, satisfies its functional
equation with root number of absolute value one, and the
uncompleted function has polynomial growth on each fixed
vertical strip.  The completion and functional equation
are recorded in \cite[\S2.1, (2.1)--(2.3)]{zaman-hecke-estimates};
the qualitative growth needed here also follows from
the convexity estimate there, Lemma~2.1, together with
the Euler product and the functional equation.
For our totally odd character the completion is
\[
 \Lambda(s)=d_F^{s/2}
       \left\{\pi^{-(s+1)/2}
                    \Gamma\left(\frac{s+1}{2}\right)\right\}^{f}
                    L_F(s,\chi).
\]
Since \(\chi\) is real, its functional equation is
\(\Lambda(s)=W\Lambda(1-s)\), where \(|W|=1\).
The Euler product and the integral test give
\[
 |L_F(3/2+it,\chi)|\le\zeta(3/2)^f\le3^f.
\]
On the other boundary of the strip, the functional equation
and \(\Gamma(z+1)=z\Gamma(z)\) therefore give
\[
 |L_F(-1/2+it,\chi)|
       \le d_F\left(\frac{3}{\pi}
                            |1/4-it/2|\right)^f.
\]
Indeed the two gamma arguments are \(5/4-it/2\) and
\(1/4+it/2\), whose quotient in absolute value is
\(|1/4-it/2|\).

Put \(F(s)=L_F(s,\chi)/(s+2)^f\).  On
\(\Re s=3/2\), \(|F(s)|\le1\); on \(\Re s=-1/2\),
\(|F(s)|\le d_F\), because
\[
 \frac{|1/4-it/2|}{|3/2+it|}\le\frac12,\qquad
                         \frac{3}{2\pi}<1.
\]
The maximum principle on this strip consequently yields
\[
 |F(\sigma+it)|\le d_F^{(3/2-\sigma)/2}.
\]
For clarity, the unbounded-strip step follows by applying
the ordinary maximum principle on rectangles to
\[
 F(s)d_F^{(s-3/2)/2}
                  \exp\{\varepsilon(s^2-9/4)\}.
\]
Its vertical boundary values have absolute value at most
one.  Polynomial vertical growth makes its horizontal
boundary values tend to zero as the rectangle grows.
Then let \(\varepsilon\) decrease to zero.  Evaluation at
\(s=1\) proves the first assertion.

By \eqref{rv-unitary-volume-formula},
\[
 \operatorname{vol}(\mathcal X)
     =2d_F^{n^2/2}\prod_{\ell=1}^n L_F(\ell,\chi^\ell).
\]
For the remaining factors the Euler products give
\[
 \prod_{\ell=2}^n |L_F(\ell,\chi^\ell)|
 \le\left(\prod_{\ell=2}^{\infty}\zeta(\ell)\right)^f
 \le e^f.
\]
The last inequality follows from
\[
 \sum_{\ell=2}^{\infty}\log\zeta(\ell)
 \le\sum_{\ell=2}^{\infty}\sum_{a=2}^{\infty}a^{-\ell}
 =\sum_{a=2}^{\infty}\frac1{a(a-1)}=1.
\]
Thus the volume is at most
\(18^f r_E^{N/2+f/4}\).
Since \(X\ge18\), \(r_E\le X\), and \(n\ge2\), this is
at most \(X^{N/2+5f/4}\le X^N\).

Finally, reduction modulo \(g\) embeds \(K_f/K_f(g)\)
in the endomorphisms of
\(\widehat\Lambda_0/g\widehat\Lambda_0\).
The latter additive group has order \(g^{2fn^2}\).
This follows locally from freeness over a discrete
valuation ring and hence holds for a projective lattice.
The index, and therefore the reciprocal level volume,
is at most \(g^{2N}\).
\end{proof}

\subsection{A small fundamental domain for a radical}

We use the least sublattice size \(\lambda(L)\) and
greatest quotient size \(\mu(L)\) of
Section~\ref{nu-section}, always with induced and quotient
positive metrics.  Conjugate duality gives
\[
            \lambda(\overline L^{\,*})=\mu(L)^{-1}.
\]
If \(\mathfrak a\) is a fractional ideal of \(E\), then
\begin{equation}\label{dc-ideal-size}
 \lambda(\mathfrak a L)
     =|N_{E/\mathbb Q}\mathfrak a|^{1/(2f)}\lambda(L),
 \qquad
 \mu(\mathfrak a L)
     =|N_{E/\mathbb Q}\mathfrak a|^{1/(2f)}\mu(L).
\end{equation}
Indeed \(M\mapsto\mathfrak aM\) is a bijection between
sublattices, and between projective quotients, in the
two modules.  On rank \(r\) covolumes it multiplies
the normalized height by
\(|N\mathfrak a|^r\), as follows by localization.
Taking \(2fr\)-th roots proves the identities.

For a real Euclidean lattice \(L\), denote its covering
radius by \(\operatorname{covrad}(L)\), its first vector
minimum by \(\lambda_1(L)\), and its Euclidean integral
dual by \(L^\vee\).  We import Banaszczyk's theorem
\begin{equation}\label{dc-transference}
       \operatorname{covrad}(L)\lambda_1(L^\vee)
                         \le\frac{\dim_{\mathbb R}L_{\mathbb R}}2.
\end{equation}
This is \cite[Theorem~2.2]{banaszczyk-transference},
also stated with these Euclidean conventions in
\cite[Theorem~3.7.2, (3.54)]{bost-kunnemann-extensions}.
It is used for arbitrary lattices, without semistability.

\begin{lemma}\label{dc-hom-cover}
Let \(A,B\) be positive metrized projective
\(\mathcal O_E\)-lattices of ranks \(a,b\ge1\).
Give \(\operatorname{Hom}_{\mathcal O_E}(A,B)\) the
Euclidean norm
\((\sum_v\|T_v\|_{\rm HS}^2)^{1/2}\).  Then
\begin{equation}\label{dc-hom-cover-bound}
 \operatorname{covrad}\operatorname{Hom}(A,B)
          \le\frac{ab\sqrt f\,r_E}{2}
                                      \frac{\mu(B)}{\lambda(A)}.
\end{equation}
\end{lemma}
\begin{proof}
Write \(\mathfrak d_E^{-1}\) for the absolute
codifferent.  The trace pairing identifies the integral
dual of \(\operatorname{Hom}(A,B)\) with
\(\mathfrak d_E^{-1}\operatorname{Hom}(B,A)\).
Under the Euclidean Riesz identification, a map \(U\)
in this module corresponds to \(2U^*\), because
\[
 \operatorname{Tr}_{E/\mathbb Q}\operatorname{tr}(UT)
                  =2\sum_v\Re\operatorname{tr}(U_vT_v).
\]
The trace pairing is perfect for projective modules:
one checks this in free local bases at every finite
prime.  Formula \eqref{dc-ideal-size} and
\(|N\mathfrak d_E^{-1}|=|D_E|^{-1}=r_E^{-2f}\)
give
\(\lambda(\mathfrak d_E^{-1}A)=r_E^{-1}\lambda(A)\).
For \(U\ne0\), Lemma~\ref{nu-map-size} and the
arithmetic--geometric mean inequality give
\[
 \left(\sum_v\|2U_v^*\|_{\rm HS}^2\right)^{1/2}
       \ge2\sqrt f\,r_E^{-1}\frac{\lambda(A)}{\mu(B)}.
\]
Apply \eqref{dc-transference} in real dimension \(2fab\).
\end{proof}

Fix a point \(x\) with \(0<\tau=\lambda_x<1/2\).
Use the canonical isotropic flag of
Proposition~\ref{nu-canonical-flag}, with successive
semistable quotient sizes
\[
             0<\tau_1<\cdots<\tau_a<1/2,\qquad
             \tau_1=\tau,\qquad a\le q.
\]
Set \(\tau_{a+1}=1/2\).

\begin{lemma}\label{dc-radical-domain}
There is a nonzero primitive isotropic sublattice \(P\)
of rank \(j\le q\) such that the following holds.
Let \(U\) be the unipotent radical of its stabilizer,
and put
\[
 \Gamma_U(g)=
 \{\gamma\in U(F):\gamma\Lambda_x=\Lambda_x,\
                         (\gamma-1)\Lambda_x\subseteq g\Lambda_x\}.
\]
In orthonormal Witt coordinates at infinity, every
class in \(\Gamma_U(g)\backslash U_\infty\) has a
representative \(\exp A\exp C\), where \(C\) is central,
\([A,C]=0\), and
\begin{equation}\label{dc-radius}
 \|A\|_{\rm HS}+\|C\|_{\rm HS}\le\rho,\qquad
 \rho=\frac{n^2\sqrt f\,g r_E}{4}(2\tau)^{1/q}.
\end{equation}
Here and below a norm on a product of real places is
the square root of the sum of the squared norms.
\end{lemma}
\begin{proof}
An extension of lattices whose greatest quotient sizes
are at most \(s\) again has greatest quotient size at
most \(s\).  To see this, take any projective quotient,
saturate the image of the sublattice, and use the exact
covolume identity between this saturated image and
the remaining quotient.  Saturation decreases
covolume; both induced maps are contractions for
quotient metrics.  Each of the two factors therefore
has size at most \(s\).  Apply this successively to the
semistable grades of \(P_i\) to obtain
\(\mu(P_i)\le\tau_i\).
The self-dual core \(P_i^\perp/P_i\) has least size
\(\tau_{i+1}\) if \(i<a\), and at least \(1/2\)
if \(i=a\), by construction of the flag.

The product of the \(a\) ratios
\(\tau_i/\tau_{i+1}\) is \(2\tau<1\).
For some \(i\),
\[
 \frac{s}{t}:=\frac{\tau_i}{\tau_{i+1}}
       \le(2\tau)^{1/a}\le(2\tau)^{1/q}=:R.
\]
Take \(P=P_i\), \(j=\operatorname{rank}P\), and
\(M=P^\perp/P\), of rank \(m=n-2j\).
Then
\[
 \mu(P)\le s,\quad
 \lambda(\overline P^{\,*})\ge s^{-1},\quad
 \lambda(M)\ge t,\quad \mu(M)\le t^{-1},\quad
                         s^2\le R/4.
\]
The last assertion uses \(s/t\le R\) and \(t\le1/2\).

Use the integral hyperbolic splitting proved in
Section~\ref{rv-section}.  At infinity use orthogonal
Witt coordinates \(P\oplus M\oplus\overline P^{\,*}\),
so the Hermitian matrix is
\[
       \begin{pmatrix}0&0&I\\0&J_0&0\\I&0&0\end{pmatrix}.
\]
The first and central Lie coordinates of the radical are
\[
 A_Z=\begin{pmatrix}0&-Z^*J_0&0\\0&0&Z\\0&0&0\end{pmatrix},
 \qquad
 C_0=\begin{pmatrix}0&0&C\\0&0&0\\0&0&0\end{pmatrix},
 \qquad C^*=-C.
\]
Thus \(\|A_Z\|_{\rm HS}=\sqrt2\|Z\|_{\rm HS}\).
Projection of \(\Gamma_U(g)\) onto the \(Z\)-coordinate
is exactly
\[
               g\operatorname{Hom}(\overline P^{\,*},M).
\]
For surjectivity, choose \(\theta\in\mathcal O_E\)
with \(\theta+\bar\theta=1\), available because
\(E/F\) is unramified at finite places.  For an integral
\(Z\) divisible by \(g\), the group matrix with top
right block \(B=-\theta Z^*J_0Z\) is integral and
congruent to the identity modulo \(g\).
This statement is intrinsic to the integral splitting;
its expression in orthonormal real coordinates need not
be integral.  It nevertheless identifies the projected
lattice and its induced metric exactly.
Lemma~\ref{dc-hom-cover} now gives first-coordinate
covering radius, in the radical norm, at most
\begin{equation}\label{dc-first-radius}
                  \frac{jm\sqrt f\,g r_E}{\sqrt2}\frac{s}{t}.
\end{equation}

We give the corresponding duality for the center,
including the prime two.  Let \(\mathcal C(P)\) be the
skew-Hermitian maps \(\overline P^{\,*}\to P\).
The \(\mathcal O_F\)-pairing
\[
 \mathcal C(P)\times\mathcal C(\overline P^{\,*})
       \longrightarrow\mathcal O_F,\qquad
                         (C,U)\longmapsto-\operatorname{tr}(CU)
\]
is perfect.  In local free bases, an off-diagonal pair
is the trace pairing for the unramified quadratic
algebra.  A diagonal coordinate is its trace-zero
rank-one module.  At an inert prime choose a generator
\(\alpha\) of the quadratic algebra with unit
discriminant; then
\(2\alpha-\operatorname{Tr}(\alpha)\) generates its
trace-zero module and its square is a unit.
This holds also at a dyadic inert prime: the trace of
such an \(\alpha\) is a unit, and the equation
\(2a+\operatorname{Tr}(\alpha)b=0\) forces
\(b\in2\mathcal O_F\).
At a split prime the generator is \((1,-1)\).
These checks prove perfectness at every finite prime,
and hence globally.

The Euclidean integral dual of \(\mathcal C(P)\) is
therefore represented by
\(\mathfrak d_F^{-1}\mathcal C(\overline P^{\,*})\).
There is no factor two here: the matrix trace already
lies in \(F\), and the real inner product is
\(-\sum_v\operatorname{tr}(C_vU_v)\).
Because
\[
 |N_{E/\mathbb Q}(\mathfrak d_F^{-1}\mathcal O_E)|
                        =d_F^{-2}=|D_E|^{-1},
\]
every nonzero dual vector, regarded as a map
\(P\to\mathfrak d_F^{-1}\overline P^{\,*}\),
has Euclidean norm at least
\(\sqrt f/(r_Es^2)\), by Lemma~\ref{nu-map-size}.
Its real dimension is \(fj^2\); transference gives
central covering radius at most
\begin{equation}\label{dc-central-radius}
                        \frac{j^2\sqrt f\,g r_E}{2}s^2.
\end{equation}

Reduce the first coordinate by an element of
\(\Gamma_U(g)\), using \eqref{dc-first-radius}, and
then reduce the remaining central coordinate by
\(g\mathcal C(P)\), using \eqref{dc-central-radius}.
Group multiplication is triangular, so the second
operation leaves the first coordinate unchanged.
An affine translation of the central fiber has the
same covering radius.  The resulting representative
is \(\exp A_Z\exp C_0\); the two matrices commute.
Its sum of norms is at most
\[
 \sqrt f\,g r_ER\left(\frac{jm}{\sqrt2}+\frac{j^2}{2}\right)
 \le\sqrt f\,g r_ER\,j(n-j)
 \le\frac{n^2\sqrt f\,g r_E}{4}R.
\]
This proves the assertion.
\end{proof}

\subsection{Zero constant terms for the actual projection}

We use the following classical Harish--Chandra theorem.
A matrix coefficient between compact-finite vectors
of a discrete-series representation of a connected
semisimple real Lie group with finite center belongs
to its Harish--Chandra cusp Schwartz space.  In particular,
for every proper parabolic radical \(U\) its integrals
\[
                   \int_{U(\mathbb R)}c(aub)\,du
\]
converge absolutely and vanish, for all fixed \(a,b\).
The Schwartz assertion is in
\cite[\S33, Lemma~65, Corollary~1]{harish-chandra-ds-ii};
the cusp assertion and its constant-term application
are explicitly recalled in
\cite[\S4.1]{khare-larsen-savin}.
We import this local theorem, rather than a statement
about the global cuspidal spectrum.

It applies to \(U(p,q)\) as follows.
The map \(U(1)\times SU(p,q)\to U(p,q)\),
\((z,g)\mapsto zg\), is onto with finite kernel.
The pullback of an irreducible representation is a
unitary central character times an irreducible
representation of \(SU(p,q)\).  Square integrability
and compact finiteness pass across this cover because
the added center is compact.  A parabolic radical
is contained in \(SU(p,q)\); its constant-term
integral is thus covered by the stated theorem.
The same reasoning applies to the contragredient
coefficient and to a finite product of real places.

\begin{lemma}\label{dc-cuspidal}
Write the projection identity of Section~\ref{ak-section}
as
\[
 K(x,y)=\sum_{i=1}^D F_i(x)\overline{F_i(y)},\qquad
 F_i=T_iw,
\]
where each \(T_i:\sigma\to L^2(\mathcal X)\) is an
isometric intertwiner and \(w\) is the fixed unit
lowest-compact-type vector.  Every \(F_i\), every
archimedean right translate of it, and every one of
its right Lie derivatives is cuspidal.
\end{lemma}
\begin{proof}
For fixed \(y\), integrate the kernel column along
\(U(F)\backslash U(\mathbb A_F)\).
This quotient is compact.  Absolute locally uniform
convergence of the kernel series allows termwise
integration, including absolute values.  Decomposing
\(G(F)\) into left \(U(F)\)-cosets unfolds this integral
to
\[
 \sum_{\gamma\in U(F)\backslash G(F)}
       \int_{U(\mathbb A_F)}
                   \Phi(x^{-1}u^{-1}\gamma y)\,du .
\]
The total integral of the absolute values is finite
by the preceding compact-quotient calculation.
Fubini is therefore applicable to the finite and
infinite factors.  Each archimedean radical integral
vanishes by the local theorem, so every kernel
column has zero constant term.

These columns span \(\operatorname{span}\{F_1,\ldots,F_D\}\).
Indeed, if their coefficient vectors did not span
\(\mathbb C^D\), a nonzero linear combination of the
\(F_i\) would vanish at every point, contradicting
their orthonormality in \(L^2\).
Consequently all \(F_i\) are cuspidal.
Right translations preserve constant terms.
For Lie derivatives, differentiation under the
constant-term integral is valid because its domain
is compact and the functions are smooth.
\end{proof}

For the particular radical of Lemma~\ref{dc-radical-domain},
cuspidality also means zero average on
\(\Gamma_U(g)\backslash U_\infty\), with \(x_f\) fixed.
Here is the finite-place justification.
If
\[
 K_{U,f}(g)=U(\mathbb A_{F,f})\cap x_fK_f(g)x_f^{-1},
\]
then
\[
                 U(\mathbb A_{F,f})=U(F)K_{U,f}(g).
\]
This follows directly from the two-step coordinates:
approximate the first additive coordinate by a rational
one, lift it by the trace-one formula above, and then
approximate the central coordinate.  Density of a
finite-dimensional \(F\)-vector space in its finite
adeles follows by clearing denominators and the
Chinese remainder theorem.  Thus the adelic
constant-term quotient, after imposing right level
\(K_{U,f}(g)\), is precisely
\(\Gamma_U(g)\backslash U_\infty\).
The first-coordinate and central lattices also show
directly that this real quotient is compact.
Normalize its Haar measure to have total mass one.

\subsection{Derivatives on a bounded range of compact types}

At one real place use the Hilbert--Schmidt inner
product on the real Lie algebra \(\mathfrak g\) of
\(U(p,q)\), and let \(\mathfrak g=\mathfrak k\oplus
\mathfrak p\) be its orthogonal Cartan decomposition.
Put
\[
                \beta(P)=2\sqrt n\,b^2+4n^{3/2}+4P.
\]
In the product representation, let \(V_P\) be the
span of all products of at most \(P\) complexified
Lie algebra operators applied to the full lowest
compact type \(W=\bigotimes_v W_v\).
It is a finite-dimensional compact-invariant space
of smooth vectors.

\begin{lemma}\label{dc-derivative}
For a Hilbert--Schmidt unit vector \(Y\) in the Lie
algebra of one real place and \(u\in V_P\),
\[
                  \|d\sigma(Y)u\|\le\beta(P)\|u\|.
\]
For a general product Lie algebra vector \(Y\),
\begin{equation}\label{dc-derivative-bound}
       \|d\sigma(Y)u\|
                   \le\sqrt f\,\beta(P)\|Y\|_{\rm HS}\|u\|.
\end{equation}
\end{lemma}
\begin{proof}
We include the Casimir computation to fix its sign.
For orthonormal bases of \(\mathfrak k,\mathfrak p\)
the invariant form is \(B(A,B)=\Re\operatorname{tr}(AB)\),
negative on \(\mathfrak k\) and positive on
\(\mathfrak p\).  Thus
\[
 \Omega_G=\sum_{\mathfrak p}Y^2-\sum_{\mathfrak k}Y^2,
 \qquad \Omega_K=-\sum_{\mathfrak k}Y^2,
 \qquad
 \sum_{\mathfrak g}\|d\sigma(Y)u\|^2
                    =\langle u,(2\Omega_K-\Omega_G)u\rangle.
\]
The Casimir eigenvalues in this normalization are
\[
 c_G=\|\lambda_v\|^2-\|\rho\|^2,\qquad
 c_K(\mu)=\|\mu+\rho_K\|^2-\|\rho_K\|^2.
\]
These formulas follow by writing the complexified
Casimir in matrix units \(E_{ij}\): move each positive
root operator to the right of its negative root
operator and apply it to a highest vector.  The
commutators contribute \(2\rho\), and the diagonal
terms contribute the squared weight.  The same
calculation for the two compact blocks gives
\(\rho_K\).  On the noncompact representation the
Harish--Chandra parameter is, by definition, the
shifted parameter for this central character.

The lowest compact highest weight is
\(\kappa=\lambda_v+\rho-2\rho_K\), as checked in
Section~\ref{fd-section}.  We use the elementary bounds
\[
 \|\rho\|,\|\rho_K\|\le\frac12 n^{3/2},\qquad
              \|\kappa\|\le\sqrt n\,b^2+\frac32n^{3/2}.
\]
All weights of the compact module of highest weight
\(\kappa\) have norm at most \(\|\kappa\|\).
For completeness, the highest-weight construction
places each weight below \(\kappa\) by a sum of
positive roots.  The set of weights is Weyl invariant.
Testing an arbitrary real linear functional, and
moving it into the dominant chamber, therefore
places every weight in the convex hull of the Weyl
orbit of \(\kappa\).  That convex hull lies in the
closed ball of radius \(\|\kappa\|\).

A highest weight \(\mu\) occurring in \(V_P\) occurs
in \(W_v\otimes\mathfrak g_{\mathbb C}^{\otimes r}\)
for some \(r\le P\).  The weights of the adjoint
module are zero or \(e_i-e_j\), of norm \(\sqrt2\).
Consequently
\[
 \|\mu\|\le\|\kappa\|+\sqrt2P,\qquad
 c_K(\mu)\le
             \bigl(\sqrt n\,b^2+2n^{3/2}+2P\bigr)^2.
\]
If the expression inside parentheses is \(a_P\),
the energy above is at most
\[
 \left(2a_P^2+\frac14 n^3\right)\|u\|^2
                                      \le4a_P^2\|u\|^2.
\]
This is \(\beta(P)^2\|u\|^2\).
The calculation applies to the compact-isotypic
decomposition of every \(u\in V_P\), including in
the product representation.  Cauchy--Schwarz for a
linear combination of the orthonormal Lie basis
proves the single-place bound.  Sum over the \(f\)
places and apply Cauchy--Schwarz once more to obtain
\eqref{dc-derivative-bound}.
\end{proof}

\subsection{An explicit local Sobolev estimate}

Write
\[
 {\bf F}(u,y)=\bigl((T_1u)(y),\ldots,(T_Du)(y)\bigr)
                         \in\mathbb C^D,\qquad
                  \tau_y=\min\{\lambda_y,1/2\}.
\]
For \(u\in V_P\) these are smooth representatives
obtained by right differentiation and compact
translation of the \(F_i\).  They satisfy pointwise
equivariance and are right \(K_f(g)\)-invariant.
Each \(T_i\) is isometric, so
\begin{equation}\label{dc-L2}
       \int_{\mathcal X}\|{\bf F}(u,y)\|^2\,dy=D\|u\|^2.
\end{equation}

\begin{lemma}\label{dc-sobolev}
For \(u\in V_P\),
\begin{equation}\label{dc-sobolev-bound}
 \|{\bf F}(u,y)\|^2
 \le
  \bigl(200N^3\beta(P+N)^2\bigr)^N
       \left(\frac7{\tau_y^2}\right)^{2N}
          \frac{D}{\operatorname{vol}(K_f(g))}\|u\|^2.
\end{equation}
\end{lemma}
\begin{proof}
Choose a Hilbert--Schmidt orthonormal basis
\(Y_1,\ldots,Y_N\) of the product real Lie algebra,
each basis vector supported at one place.  Set
\[
 r=\frac1{20N^2},\qquad
 \Psi(t)=\exp(t_1Y_1)\cdots\exp(t_NY_N),
                         \quad t\in[-r,r]^N.
\]
We record explicit chart bounds.  Put \(S=\sum_j|t_j|\),
so \(S\le1/(20N)\).  In the ambient product of
matrix spaces, the \(j\)-th derivative column differs
from \(Y_j\) by norm at most \(e^S-1\le2S\).
The operator norm of the difference between the full
derivative and the linear isometry
\(t\mapsto\sum t_jY_j\) is at most
\(2S\sqrt N\le1/10\).
Integrating along a straight segment in the cube
shows that \(\Psi\) is injective.
Its left Maurer--Cartan columns are
\[
 A_j(t)=\operatorname{Ad}
       \bigl(\exp(-t_NY_N)\cdots\exp(-t_{j+1}Y_{j+1})\bigr)Y_j.
\]
They differ from \(Y_j\) by norm at most
\(e^{2S}-1\le3S\).  Their full derivative matrix
has smallest singular value at least
\(1-3S\sqrt N>1/2\).
Thus Haar measure in the chart has density at least
\(2^{-N}\), and \(\|A_j(t)\|_{\rm HS}<2\).

The map of the chart times \(K_f(g)\) into
\(\mathcal X\), translated by \(y\), has multiplicity
at most
\[
                    M_y=\left(\frac7{\tau_y^2}\right)^{2N}.
\]
Indeed any two lifts differ by a rational isometry
preserving \(\Lambda_y\), of operator norm less than
\(e^{2Nr}<2\) at every real place.
The difference of two such isometries is a nonzero
integral endomorphism.  Lemma~\ref{amb-map-separation}
gives geometric mean operator norm at least
\(\lambda_y^2\).  The packing argument of
Lemma~\ref{amb-packing}, in real matrix dimension
\(2n^2\) at each place, bounds their number by
\[
                  (1+4/\tau_y^2)^{2N}\le M_y.
\]
Consequently, for any fixed smooth vector \(v\)
in the span under consideration,
\begin{equation}\label{dc-chart-integral}
 \int_{[-r,r]^N}\|{\bf F}(v,y\Psi(t))\|^2\,dt
       \le\frac{2^N M_yD}{\operatorname{vol}(K_f(g))}\|v\|^2.
\end{equation}
This follows from the chart Jacobian, the multiplicity
bound, right level invariance, and \eqref{dc-L2}.

Here is an elementary Sobolev inequality sufficient
for this chart.  The one-dimensional fundamental
theorem of calculus, averaged over \([-r,r]\),
expresses \(h(0)\) as its average plus an integral
of \(h'\) with kernel of absolute value at most
\(1/2\).  Iterating in the \(N\) coordinates gives
\[
 h(0)=\sum_{I\subseteq\{1,\ldots,N\}}
          \int_{[-r,r]^N}k_I(t)\partial_Ih(t)\,dt,
 \qquad |k_I(t)|\le(2r)^{-N}r^{|I|}.
\]
The same identity holds for a finite-dimensional
Hilbert-space-valued function.  Cauchy--Schwarz gives
\[
 \|h(0)\|^2\le(2r)^{-N}(1+r^2)^N
                     \sum_I\int\|\partial_Ih(t)\|^2\,dt
              \le r^{-N}\sum_I\int\|\partial_Ih(t)\|^2\,dt.
\]
Use \(h(t)={\bf F}(u,y\Psi(t))\).
For \(I=\{i_1<\cdots<i_k\}\), direct differentiation
of the ordered product gives
\[
 \partial_Ih(t)=
 {\bf F}\bigl(d\sigma(A_{i_1}(t))\cdots
                       d\sigma(A_{i_k}(t))u,\ y\Psi(t)\bigr).
\]
Each coordinate is differentiated only once, so this
identity has no additional derivative terms.
Expand the \(A_j(t)\) in the fixed orthonormal basis.
The squared sum of coefficients of their tensor
product is at most \(4^k\); pointwise Cauchy--Schwarz
therefore bounds the squared norm by \(4^k\) times
the sum for all \(N^k\) fixed words of length \(k\).
Lemma~\ref{dc-derivative} bounds each such word by
\(\beta(P+N)^k\|u\|\).
Applying \eqref{dc-chart-integral} word by word yields
\[
 \|{\bf F}(u,y)\|^2
 \le r^{-N}\frac{2^NM_yD}{\operatorname{vol}(K_f(g))}
                 (1+4N\beta(P+N)^2)^N\|u\|^2.
\]
Since \(N\ge1\), \(\beta(P+N)\ge1\), and
\(r^{-1}=20N^2\), this proves the stated constant.
\end{proof}

\subsection{Iteration of the zero-average identity}

Fix \(x\) with \(\tau=\lambda_x<\tau_0\), and use the
radical and \(\rho\) of Lemma~\ref{dc-radical-domain}.
Let \(\Omega_x\) consist of the points
\(x\exp A\exp C\) with first and central radical
coordinates satisfying
\(\|A\|_{\rm HS}+\|C\|_{\rm HS}\le\rho\);
the finite component remains \(x_f\).
It contains \(x\), the paths
\(x\exp(t(A+C))\), \(0\le t\le1\), and representatives,
up to left rational and right \(K_f(g)\) action,
of all points on the radical quotient through \(x\).

If \(\rho\le\log2\), then throughout this set
\begin{equation}\label{dc-size-nearby}
                   \tau_y\ge\tau/2.
\end{equation}
Indeed at place \(v\) the positive metrics change
by a map and inverse of norm at most
\(\exp(\|A_v\|_{\rm HS}+\|C_v\|_{\rm HS})\).
Taking the normalized product over the \(f\) places
changes every sublattice size by at most
\(\exp(\rho/\sqrt f)\).  Since \(\tau<1/2\),
the same lower bound survives truncation at \(1/2\).

Set
\[
 P_0=8qN,\qquad \beta_*=\beta(P_0+N).
\]
Our parameters give the following simple bounds:
\begin{equation}\label{dc-numerical}
 \beta_*\le X^3,\qquad N\le X^2,\qquad
 156800\le X,\qquad
 \rho\le\tfrac12X^{-9}<\log2.
\end{equation}
For the first, \(q\le\sqrt n\) implies
\[
 \beta_*\le2\sqrt n\,b^2+4n^{3/2}+36n^{5/2}f.
\]
Divide by \(n^3f^3b^3\le X^3\); the quotient is at most
\[
        n^{-5/2}+\tfrac12n^{-3/2}+\tfrac92n^{-1/2}<1
                         \qquad(n\ge100).
\]
The other elementary bounds follow from \(X\ge2n\),
\(\log n\ge20\), and
\[
 \rho\le\frac{n^2\sqrt f\,gr_E}{2}\tau^{1/q}
                          \le\tfrac12X^2X^{-11}.
\]
Equations \eqref{dc-sobolev-bound},
\eqref{dc-size-nearby}, and \eqref{dc-numerical}
give, uniformly for \(y\in\Omega_x\) and \(u\in V_{P_0}\),
\begin{equation}\label{dc-initial-evaluation}
 \|{\bf F}(u,y)\|^2
   \le X^{13N}\tau^{-4N}
       \frac{D}{\operatorname{vol}(K_f(g))}\|u\|^2.
\end{equation}
Here \(200\cdot28^2=156800\):
the Sobolev constant contributes \(N^3\le X^6\),
\(\beta_*^2\le X^6\), and the remaining constant
contributes at most \(X\).

\begin{lemma}\label{dc-iteration}
Suppose \(M_{P+1}\) bounds
\(\|{\bf F}(u,y)\|/\|u\|\) for all nonzero
\(u\in V_{P+1}\), \(y\in\Omega_x\).
For \(0\le P<P_0\), the same quantity on \(V_P\)
is bounded by
\[
                   2\rho\sqrt f\,\beta_* M_{P+1}.
\]
\end{lemma}
\begin{proof}
For \(y=x\exp A\exp C\in\Omega_x\), differentiate
along \(x\exp(t(A+C))\).  The two coordinates commute,
and \(d\sigma(A+C)u\in V_{P+1}\).
Lemma~\ref{dc-derivative} gives
\[
 \|{\bf F}(u,y)-{\bf F}(u,x)\|
                  \le\rho\sqrt f\,\beta_*M_{P+1}\|u\|.
\]
The same inequality holds for every point of the
radical quotient, by its representatives in
\(\Omega_x\) and the rational and level invariance.
Its average is zero by Lemma~\ref{dc-cuspidal} and
the finite-place argument following it.
Averaging the difference therefore bounds
\(\|{\bf F}(u,x)\|\) by the same right side.
The triangle inequality proves the assertion at
all \(y\in\Omega_x\).
\end{proof}

\begin{proof}[Proof of Theorem~\ref{dc-main}]
If \(D=0\), the finite-rank identity makes the kernel
zero, and there is nothing to prove.
Otherwise iterate Lemma~\ref{dc-iteration} \(P_0\)
times, starting with \eqref{dc-initial-evaluation}.
Since \(w\in V_0\) has norm one, the diagonal identity
gives
\[
 K(x,x)\le(2\rho\sqrt f\,\beta_*)^{16qN}
       X^{13N}\tau^{-4N}
                   \frac{D}{\operatorname{vol}(K_f(g))}.
\]
Moreover
\[
 2\rho\sqrt f\,\beta_*
   =\frac{n^2fgr_E}{2}\,2^{1/q}\beta_*\tau^{1/q}
                         \le X^5\tau^{1/q},
\]
because \(q\ge1\), \(n^2fgr_E\le X^2\),
and \(\beta_*\le X^3\).
It follows that
\[
 K(x,x)\le
 X^{(80q+13)N}\tau^{12N}
                   \frac{D}{\operatorname{vol}(K_f(g))}
 \le X^{(-52q+13)N}
                   \frac{D}{\operatorname{vol}(K_f(g))}.
\]
The diagonal is nonnegative.  Integrating this
uniform bound over the deep subset, bounding its
volume by the whole quotient volume, and using
Lemma~\ref{dc-volume}, gives
\[
 \int_{\lambda_x<\tau_0}K(x,x)\,dx
 \le X^{(-52q+16)N}D
 \le X^{-36N}D
 \le e^{-N}D.
\]
This is \eqref{dc-main-bound}.
\end{proof}

\begin{corollary}\label{dc-ambient-comparison}
Assume the field, lattice, finite test-function,
congruence, and discrete-series hypotheses in
Section~\ref{amb-section}.  Suppose \(C_g\ge0\),
\(b\ge2\), \(1\le q\), all the quantitative conditions
\eqref{amb-parameters} hold, and
\[
              \log n\ge20,\qquad n\ge4\log(C_g+3).
\]
Then the actual ambient multiplicity is finite and
\[
             |D-D_*|\le
                    \exp(-nf\log n/16)\,D_*.
\]
\end{corollary}
\begin{proof}
The compact-type estimate of Section~\ref{amb-section}
supplies the local decay required in
Section~\ref{ak-section}.  The lower bound for \(h\)
in \eqref{amb-parameters} in particular implies its
condition \(h/(8\sqrt{2q})>2n^2\).
That section therefore constructs the finite-rank
orthogonal projection, its absolutely convergent
continuous kernel, and its pointwise diagonal.
Theorem~\ref{nu-main} proves \eqref{amb-near-input}.
The rare-cusp estimate \eqref{rv-probability} implies
\eqref{amb-rare-input}, since
\(\tau_0<1/2\) and \(e^{-nf\log n/3}\le e^{-nf\log n/4}\).
Theorem~\ref{dc-main} proves \eqref{amb-deep-input},
and Theorem~\ref{amb-far} supplies the distant terms.
All hypotheses of Proposition~\ref{amb-final-reduction}
are now proved, so it yields the asserted comparison.
\end{proof}

This corollary concerns the complex ambient multiplicity.
It retains the explicitly stated classical imports of
the preceding sections.  It does not assert that
the remaining arithmetic and functorial steps of
the submitted GRH argument have been verified.


